\scanpage{75}
On s'attend donc à pouvoir écrire % 1
$A(S)\otimes_{\mathbb{Z}}\mathbb{Q}$ en termes de ce motif $M$. % 2
Or en effet, des \unsure{considérations|considér.|raisons} heuristiques convaincantes % 3
\deleted{pour} nous suggèrent % 4
\[
 A(S)\otimes_{\mathbb{Z}}\mathbb{Q}\simeq H^1(S,T_{\mathrm{mot}}(A)) % 5
\]
[NB. On doit avoir $H^0(S,T_{\mathrm{mot}}(A))=0$, % 6
à cause de \deleted{Mordell-Weil} considérations de % 7
poids ---- ]. \deleted{Une autre façon de} \inserted{Considérons} (pour $\ell$ premier % 8
aux car. résiduelles \unsure{de|à} $X$) % 9
\[
 \unsure{\varnothing|0}\longrightarrow{}_{\ell^\nu}A\longrightarrow A\longrightarrow A\;-\;0 % 10
\]
d'où % 11
\[
 0\longrightarrow H^0(S,A)_{\ell^\nu}\longrightarrow H^1(S,{}_{\ell^\nu}A)
 \longrightarrow{}_{\ell^\nu}H^1(S,A)\;\unsure{-|\to}\;0 % 12
\]
et par passage à la limite et $\otimes_{\mathbb{Z}_\ell}\mathbb{Q}_\ell$ % 13
\[
\begin{aligned}
0\longrightarrow {}&
\underbrace{H^0(S,A)}_{\substack{\text{de type fini}\\\text{sur }\mathbb{Z}}}
\otimes_{\mathbb{Z}}\deleted{\unsure{\mathbb{Z}|\mathbb{Q}}}\mathbb{Q}_\ell
 \longrightarrow H^1(S,T_{\ell}^{\deleted{\unsure{\mathbb{Q}}}}(A))_{\mathbb{Q}_\ell}
 \longrightarrow T_{\unsure{\bullet|\ell}}^{\deleted{\unsure{\mathbb{Q}}}}(H^1(S,A))_{\mathbb{Q}_\ell}\to0\\ % 14
&\hspace{36mm}\rotatebox{90}{$\simeq$}\\ % 15
&\hspace{32mm}H^1(S,T_{\mathrm{mot}}(A))\otimes_{\mathbb{Q}}\mathbb{Q}_{\ell}
\end{aligned} % 16
\]
et le résultat escompté ici équivaut % 17
moralement à celui-ci : % 18
\[
\boxed{\begin{aligned}
H^1(S,A)(\ell)\quad\text{\uuline{fini}}\qquad&\text{si }\ell\ \text{\textquotedblleft premier à } \deleted{\unsure{X}}S\text{\textquotedblright}\\ % 19
&\text{(}S\text{ de type fini sur }\operatorname{Spec}\mathbb{Z}\text{)}
\end{aligned}} % 20
\]
p.ex. si $S=\operatorname{Spec}\mathbb{F}_q$, alors c'est vrai par Lang ! % 21
% archive sheet number: 74
